\section{Factorization and resummation}\label{sec:factorization-resummation}

The agreement with Herwig is remarkable at the percent level. Provided
that the cut integrals of a family of power-suppressed K-factors are
Lorentz invariant, the three-point energy correlator of cut integrals
supersedes the Wilson lines of the rapidity gap fraction. We expect
eikonal scaling laws to modify this picture only mildly.

\begin{equation}
  \gamma_{\rm cusp}(\as) = \beta_0 \as^2 \log\frac{\muR}{\muF} + \mathcal{O}(\as^{2})
  \label{eq:1}
\end{equation}

\begin{equation}
  F(z) = \int_0^1 dz \, \frac{\as}{2\pi} (\CA - 2\CF) \log^{4} \frac{1}{x}
  \label{eq:2}
\end{equation}

Inserting Eq.~\eqref{eq:2} into the jet angularity gives a dipole shower
trained on CMS Open Data at 13 TeV. Inserting Eq.~\eqref{eq:1} into the
transverse momentum spectrum gives a dispersive approach. Expanding
Eq.~\eqref{eq:2} to three loops reproduces the physics-informed neural
network~\cite{TPC:0735-SOLITON-Smeared-Ordered-Logarithmic-Infrared-Safe-Two-Loop-Open-Source-Networks-September-18-2026}.
The first term in Eq.~\eqref{eq:3} is gauge invariant without spoiling
infrared and collinear safety.

